\documentclass[../../../main.tex]{subfiles}

\begin{document}
% Exercise 1.4:3(b)

\begin{proof}
	Suppose that $m < n$.
	Then $n - m$ is positive.
	by part (a),
	\[ 
		n_F - m_F = (n - m)_F = (\underbrace{1 + 1 + \ldots + 1}_{n - m \text{ times}})_F = \underbrace{1_F + 1_F + \ldots + 1_F}_{n - m \text{ times}}.
	 \]
	Since $1_F > 0_F$ by Rudin's proposition 1.18(d),
	it follows that $n_F - m_F$ is positive, Hence $m_F < n_F$.

	Conversely, suppose that $m_F < n_F$.
	Then $0 < n_F - m_F$.
	We first observe that $n - m \neq 0$, since $n - m = 0$ would imply
	\[ 
		n_F - m_F = (n - m)_F = 0_F,
	 \]
	a contradiction.
	Moreover, $n - m$ cannot be negative. Indeed, if $n - m < 0$, then
	\[ 
		n_F - m_F = (n - m)_F = -(m - n)_F,
	 \]
	where $(m - n)_F$ is positive. Thus $n_F - m_F$ would be negative by Rudin's proposition 1.18(a), again a contradiction.
	Therefore, $n - m > 0$, and hence $m < n$.

	Finally, suppose that $n \neq m$. Without loss of generality, assume that $n < m$.
	By what we have just shown, $n_F < m_F$ and therefore $n_F \neq m_F$.
	Hence the map $n \mapsto n_F$ is one-to-one.

\end{proof}

\end{document}
